\section*{Bab \ref{ch:b3:conservation-biology} --- Biologi Konservasi}

\begin{solution}{exo:b3:conservation-biology:1}
\omterm{def:b3:conservation-biology:small}{Stokastisitas demografi}: kebetulan pada kelahiran, kematian dan kelamin
beberapa individu --- kakapo pada lima puluh satu ekor, dengan jantan
lebih banyak daripada betina. \omterm{def:b3:conservation-biology:small}{Stokastisitas lingkungan}: sebuah tahun buruk
yang menerpa tiap individu sekaligus --- kekeringan, penyakit, musim
dingin yang keras yang menimpa populasi beberapa ratus ekor. \omterm{def:b3:conservation-biology:small}{Efek Allee}:
pertumbuhan per kapita yang jatuh pada kerapatan rendah --- merpati
penumpang, yang koloninya tidak dapat berbiak begitu menipis, dan yang
pemangsa tersisanya tidak lagi dijenuhkan.
\end{solution}

\begin{solution}{exo:b3:conservation-biology:2}
Ukuran sebuah populasi ideal --- tetap, berkawin acak, kelamin seimbang,
ukuran keluarga Poisson --- yang akan kehilangan keheterozigotan pada laju
yang sama dengan populasi nyatanya. Cacah jantan dan betina pembiak yang
timpang; ayunan ukuran, yang membobot generasi terkecilnya; ragam ukuran
keluarga di atas Poisson, dengan beberapa individu menghasilkan sebagian
besar anaknya (juga perkawinan tak acak dan generasi yang bertumpang
tindih).
\end{solution}

\begin{solution}{exo:b3:conservation-biology:3}
$S = cA^{z}$. Pecahan yang ditahan $= (A'/A)^{z}$: $0.5^{0.25} = 0.84$;
$0.1^{0.25} = 0.56$; $0.01^{0.25} = 0.32$ --- seperseratus luasnya masih
menahan sepertiga spesiesnya, pada akhirnya.
\end{solution}

\begin{solution}{exo:b3:conservation-biology:4}
$N_{e} \ge 50$ menahan kenaikan silang dalam di bawah $1/100$ tiap
generasi, yaitu laju yang membuat peternak menganggap depresinya masih
dapat ditoleransi: ia menjaga terhadap hilangnya kebugaran jangka pendek. $N_{e} \ge
500$ mengimbangi kehilangan keragaman oleh hanyutan terhadap pasokannya
oleh mutasi, sehingga populasinya menahan keragaman kuantitatif yang
diperlukannya untuk menyesuaikan diri: ia menjaga jangka panjangnya.
Populasi nyata memerlukan lima sampai sepuluh kali angka ini pada cacah jiwanya.
\end{solution}

\begin{solution}{exo:b3:conservation-biology:5}
$N_{e} = 4\times 8\times 120/128 = 30$; keheterozigotannya turun sebesar $1/60 =
\qty{1.7}{\%}$ tiap generasi; seperempatnya hilang ketika $(1 - 1/60)^{t} =
0.75$, $t = \ln 0.75/\ln(59/60) = 17$ generasi.
\end{solution}

\begin{solution}{exo:b3:conservation-biology:6}
$1/N_{e} = \tfrac{1}{5}(4/1000 + 1/20) = 0.0108$, $N_{e} = 93$: satu
generasi pada dua puluh telah membuat lima generasi berlaku seperti sembilan puluh tiga.
Keheterozigotan yang ditahan $(1 - 1/186)^{5} = 0.973$, berbanding $(1 -
1/2000)^{5} = 0.9975$ bagi seribu yang tetap.
\end{solution}

\begin{solution}{exo:b3:conservation-biology:7}
$F = 1 - (1 - 1/20)^{5} = 0.23$; kebugaran turun $5\times 2.3 = \qty{11}{\%}$.
Tiap anak dari betina Texas dan jantan Florida memiliki $F = 0$; bila
delapan pendatangnya adalah dua perlima betina pembiaknya, dua perlima
generasi berikutnya bersilang luar dan rata-rata $F$ jatuh menjadi sekitar
$0.6\times 0.23 = 0.14$ dalam satu generasi --- dan itulah yang ditunjukkan
pemulihannya.
\end{solution}

\begin{solution}{exo:b3:conservation-biology:8}
$r = \ln(100/31)/8 = \qty{0.15}{yr^{-1}}$. Bila berlanjut sampai 2020:
$100\,\mathrm{e}^{0.15\times 17} \approx 1200$. Tamannya menampung sekitar
seratus ekor: wilayahnya terisi, rusa wapitinya merosot, kawanannya saling
membunuh anggota, lalu kudis dan distemper menyebar --- kebergantungan
kerapatan mulai berlaku begitu jangkauan kosongnya terhuni.
\end{solution}

\begin{solution}{exo:b3:conservation-biology:9}
(a) Kemerosotan \qty{60}{\%} dalam sepuluh tahun melampaui \qty{50}{\%}: Terancam.
(b) $180 < 250$ individu dewasa: Terancam Kritis. (c) Sebarannya
\qty{3000}{km^2}, di bawah \qty{5000}{}: Terancam. (d) \num{30000}
individu dan kemerosotan \qty{10}{\%}: tidak satu pun ambangnya terpenuhi
--- tidak terancam (paling jauh Hampir Terancam, yang harus diawasi).
\end{solution}

\begin{solution}{exo:b3:conservation-biology:10}
Semua $N$ pasangan gagal dengan peluang $(1/2)^{N}$: $0.25$, $0.031$,
$0.001$, $10^{-6}$. Risiko demografis jatuh secara eksponensial bersama
$N$ lalu dapat diabaikan di luar beberapa lusin individu; di atas itu,
risiko yang tersisa bersifat lingkungan, yang tidak jatuh secepat itu,
karena ia menerpa semua individu sekaligus.
\end{solution}

\begin{solution}{exo:b3:conservation-biology:11}
Satu serpihan sepersepuluh: $0.1^{0.25} = 0.56$ spesies asalnya. Sepuluh
serpihan menampung antara $0.56$ (bila semuanya menampung spesies yang
sama) dan, pada asasnya, lebih banyak daripada hutan utuhnya (bila
masing-masing menampung spesies yang berbeda --- mustahil melampaui
lungkang kawasannya). Di mana letaknya bergantung pada pergantian antar
serpihan: habitat yang berbeda pada serpihan yang berbeda mendorong
totalnya naik; spesies yang memerlukan luas yang besar, habitat bagian
dalam, atau yang hilang dari tiap serpihan oleh \omterm{prop:b3:conservation-biology:area}{efek tepi} yang sama
mendorongnya turun; dan utang kepunahan tiap serpihan dibayar terpisah,
sehingga totalnya dalam jangka panjang biasanya di bawah total hutan utuhnya.
\end{solution}

\begin{solution}{exo:b3:conservation-biology:12}
Hanyutan menyingkirkan keheterozigotan pada laju $1/2N_{e}$ tiap
generasi; pendatang mengganti sebuah pecahan $m$ lungkang gennya dengan
alel yang diambil dari tempat lain. Dengan $mN_{e} \approx 1$, $m \approx 1/N_{e}$, yaitu dua kali
laju kehilangannya, sehingga keragamannya diisi ulang lebih cepat daripada
peluruhannya dan populasinya tetap dekat dengan sumbernya pada frekuensi
alel ($F_{ST} = 1/(1
+ 4N_{e}m) \approx 0.2$). Ongkosnya: alel yang disukai setempat diencerkan
sebesar $m$ tiap generasi, sehingga penyesuaian setempat bertahan hanya
bagi sifat yang koefisien seleksinya melampaui kira-kira $1/N_{e}$;
perbedaan setempat yang lemah terseleksi akan tertenggelamkan.
\end{solution}

\begin{solution}{pb:b3:conservation-biology:1}
\textbf{1.} $N_{e} = 4\times 20\times 40/60 = 53$, berbanding cacah jiwa 60.
\textbf{2.} $1/N_{e} = \tfrac{1}{5}(1/500 + 1/100 + 1/60 + 1/51 + 1/60) =
0.0130$, $N_{e} = 77$; generasi leher botol berisi 51 dan 60 ekor
menguasainya, yang 500 ekor nyaris tidak berarti.
\textbf{3.} $F = 1 - (1 - 1/154)^{5} = 0.032$; linier $5/154 = 0.032$.
\textbf{4.} $6\times 0.32 = \qty{1.9}{\%}$.
\textbf{5.} $(1 - 1/154)^{5} = 0.968$: \qty{96.8}{\%} ditahan.
\textbf{6.} Sepuluh generasi pada $N_{e} = 53$: tambahannya $1 - (1 -
1/107)^{10} = 0.090$, total $F = 1 - 0.968\times 0.910 = 0.12$; kehilangan
kebugarannya $6\times 1.2 = \qty{7}{\%}$; satu abad.
\textbf{7.} Pasangan asli--asli adalah $(60/70)^{2} = 0.73$ perkawinannya
dan anaknya menahan $F \approx 0.03$ (ditambah sebuah tambahan kecil);
semua yang lain memiliki $F = 0$: rata-rata $F \approx 0.73\times 0.035 = 0.026$, yaitu sebuah penurunan
seperempat dalam satu generasi.
\textbf{8.} $F$ adalah keserupaan menurut keturunan di dalam seorang
individu, dan tiap anak silang luar menyetelnya kembali ke nol;
keheterozigotan adalah sifat frekuensi alel, dan alel para migran menyebar
hanya seiring diturunkannya. Di dalam populasi yang masih kecil,
hanyutannya akan menyingkirkan alel itu lagi dalam puluhan generasi,
sehingga penyelamatannya harus diulang atau populasinya ditumbuhkan.
\textbf{9.} $N_{e} = N$ bagi kelamin yang seimbang: 500 ekor dewasa. Dari 70 ekor pada $r =
0.05$: $\ln(500/70)/0.05 = 39$ tahun.
\textbf{10.} Populasi idealnya memiliki ukuran keluarga Poisson, dengan
ragam sama dengan rata-ratanya yang dua; $N_{e} = 4N/(2 + V_{k})$,
sehingga membuat tiap pasangan menyumbang sama banyak ($V_{k} = 0$)
memberikan $N_{e} = 2N$: tidak ada galur yang hilang oleh untung-untungan sarangnya.
\textbf{11.} Kumpulkan dan simpanlah maninya untuk inseminasi buatan, lalu
pakailah pada beberapa betina tak berkerabat sepanjang beberapa tahun ---
ongkosnya adalah bahwa alelnya dapat menguasai sebuah lungkang yang kecil,
sehingga jatahnya harus dibatasi mendekati $1/N_{e}$; atau awetkanlah
jaringannya secara beku untuk pengklonan atau penghasilan gamet kelak,
dengan ongkos berupa penundaan dan risiko teknis. Tidak berbuat apa-apa
akan menghilangkan alel galur itu untuk selamanya.
\textbf{12.} \qty{160}{} juta dolar bagi sekitar 430 ekor burung:
kira-kira \num{370000} tiap ekor. Mahal tiap burung nuri, murah terhadap
sebuah spesies --- dan kurang daripada beberapa kilometer jalan tol, yaitu
perbandingan yang sungguh dibuat oleh anggaran.
\textbf{13.} $(300/2000)^{0.25} = 0.62$: sekitar 112 spesies dari 180
ditahan, dengan utang kepunahan sebesar 68.
\textbf{14.} Tiap serpihan $(100/2000)^{0.25}\times 180 = 85$ spesies.
Seluruhnya berbeda: sampai 255, mustahil di atas lungkang 180 spesies;
serupa: 85. Yang menentukan adalah seberapa berbeda habitat serpihannya
dan berapa banyak spesies yang dapat lestari dalam \qty{100}{km^2}; secara
nyata antara 85 dan 112.
\textbf{15.} $(320/2000)^{0.25}\times 180 = 114$, berbanding 85 bagi tiga
serpihan serupa yang terpencil: koridornya membiarkan penghunian kembali
dan aliran gen membuat tiga populasi kecil berlaku sebagai satu populasi yang lebih besar.
\textbf{16.} $N_{e} = 50$ dengan kelamin seimbang berarti 25 pasangan: \qty{1250}{km^2}.
Seluruh \qty{300}{km^2} menampung enam pasangan, $N_{e} \approx 12$;
sebuah serpihan, dua pasangan. Tidak satu pun mencukupi: pemangsanya
hilang berapa pun jumlah spesiesnya, dan melindungi luas yang
diperlukannya melindungi semua yang lebih kecil --- sebuah spesies payung.
\textbf{17.} $\ln(200/10)/0.03 = 100$ tahun: utangnya dibayar sepanjang
satu abad, lama setelah hutannya ditebang.
\textbf{18.} Tepinya melucuti habitat bagian dalam dari tiap serpihan,
sehingga luas yang dapat dipakai lebih kecil daripada luas yang dipetakan
dan kehilangannya lebih besar daripada yang diramalkan hukumnya; sebuah
matriks pertumbuhan sekunder atau pagar tanaman yang dapat dipakai
beberapa spesies menambahkan luas efektif lalu memperkecil kehilangannya.
\textbf{19.} $r = \ln(100/31)/8 = \qty{0.146}{yr^{-1}}$; waktu lipat duanya
$\ln 2/0.146 = 4.7$ tahun.
\textbf{20.} $N(8) = 120/[1 + (120/31 - 1)\mathrm{e}^{-1.17}] = 120/(1 +
2.87\times 0.31) = 63$. Angka 100 yang teramati adalah nilai eksponensial
itu sendiri --- $r$ dipaskan justru pada kedua titik itu --- berhadapan
dengan 63 yang logistik: pada tahun-tahun pertamanya serigalanya tidak
menemui kebergantungan kerapatan --- wilayah yang kosong dan rusa wapiti
yang berlimpah.
\textbf{21.} $\ln(4000/\num{17000})/20 = -0.072$: \qty{7}{\%} setahun.
Seratus ekor serigala yang mengambil \num{2000} ekor rusa wapiti setahun
dari kawanan yang rata-rata \num{8000} ekor berarti seperempatnya tiap
tahun --- lebih dari cukup untuk menggerakkan kemerosotannya sendiri,
meskipun perburuan di luar tamannya, kekeringan dan pemangsaan anak rusa
oleh beruang cokelat turut menyumbangnya.
\textbf{22.} Enam puluh individu dewasa, di bawah 250: Terancam Kritis.
Pemangsanya pada $N_{e} \approx 12$ (beberapa lusin ekor hewan): Terancam
Kritis pula.
\textbf{23.} Dari $3\times 10^{9}$ menjadi $1$ dalam 44 tahun: $r = -\ln(3\times
10^{9})/44 = -0.50$ tiap tahun, yang membelah dua tiap tujuh belas bulan.
Tidak ada panen yang menyingkirkan separuh dari miliaran tiap tahun;
jatuhnya lebih lambat pada mulanya dan jauh lebih cepat pada akhirnya,
karena seekor pembiak berkoloni di bawah kerapatan ambang berhenti berbiak
sementara pemangsanya tidak lagi dijenuhkan --- \omterm{def:b3:conservation-biology:small}{efek Allee} mengubah
kemerosotan menjadi keruntuhan.
\textbf{24.} Pada 1850 populasinya berjumlah miliaran, laju pertumbuhan
yang terukurnya positif dan ragamnya kecil, dan sebuah analisis kelayakan
memproyeksikan dinamika yang diberikan kepadanya: ia akan menaruh
risikonya pada nol. Ia akan melewatkan paksaan dari luar --- kereta api
dan telegraf yang membuat pemburu dapat mengikuti lalu mengosongkan tiap
tempat bertenggernya, penebangan hutan --- dan ambang Allee pada seekor
pembiak berkoloni, yang keduanya tidak ada di dalam demografi sebuah
spesies yang berlimpah; sebuah PVA melakukan ekstrapolasi, dan penyebab
kepunahan biasanya baru.
\textbf{25.} $N_{e} \approx 53$ hari ini; $F \approx 0.12$ setelah sepuluh
generasi lagi pada 60 ekor (\qty{7}{\%} kebugarannya); hutan yang menyusut
itu menahan kira-kira \qty{62}{\%} spesiesnya.
\end{solution}
